%-----------------------------------------------------------------------
% Appendix C: API Documentation
%-----------------------------------------------------------------------

\label{app:c}

\section{Introduction}
\label{app:c:introduction}

This appendix documents the APIs and interfaces of the HPC Sorting Serious Game, focusing on the signaling server REST/SSE endpoints and the key GDScript classes. The signaling server is the only custom server-side component; the game client is built in Godot Engine 4.x with GDScript.

%-----------------------------------------------------------------------
\section{Signaling Server API}
\label{app:c:signaling-server-api}

The signaling server is a Python application built on \texttt{aiohttp}. It serves three route groups: HTTP session management, HTTP lobby management with SSE, and WebSocket signaling for WebRTC.

\subsection{Session Endpoints}
\label{app:c:session-endpoints}

These endpoints manage rooms for WebRTC signaling.

\begin{table}[htbp]
    \centering
    \caption{Session management endpoints}
    \begin{tabular}{@{}llp{7cm}@{}}
        \toprule
        \textbf{Method} & \textbf{Path}                      & \textbf{Description}                                                                                                   \\
        \midrule
        \texttt{POST}   & \texttt{/session/host}             & Create a new room. Body: \texttt{\{channel, lobby\_name, public, player\_limit\}}. Returns \texttt{\{code, ws\_url\}}. \\
        \texttt{POST}   & \texttt{/session/join/\{code\}}    & Join a room by code or lobby name. Returns \texttt{\{code, ws\_url\}}.                                                 \\
        \texttt{POST}   & \texttt{/session/update/\{code\}}  & Update room metadata (lobby name, public flag, player limit).                                                          \\
        \texttt{POST}   & \texttt{/session/players/\{code\}} & Update player count for a room.                                                                                        \\
        \texttt{POST}   & \texttt{/session/close/\{code\}}   & Close a room and its associated lobby.                                                                                 \\
        \bottomrule
    \end{tabular}
\end{table}

\subsection{HTTP Lobby Endpoints (SSE)}
\label{app:c:lobby-endpoints}

These endpoints replace WebSocket-based lobby handling for browser compatibility. Real-time events are delivered via Server-Sent Events (SSE).

\begin{table}[htbp]
    \centering
    \caption{HTTP lobby endpoints}
    \begin{tabular}{@{}llp{7cm}@{}}
        \toprule
        \textbf{Method} & \textbf{Path}                  & \textbf{Description}                                                                                         \\
        \midrule
        \texttt{POST}   & \texttt{/api/lobby/connect}    & Connect and receive a peer ID. Optional body: \texttt{\{client\_id\}} to specify a GDSync client ID.         \\
        \texttt{POST}   & \texttt{/api/lobby/disconnect} & Disconnect from the server.                                                                                  \\
        \texttt{POST}   & \texttt{/api/lobby/create}     & Create a new lobby. Body: \texttt{\{lobby\_name, public, player\_limit\}}. Returns \texttt{\{code, lobby\}}. \\
        \texttt{POST}   & \texttt{/api/lobby/join}       & Join a lobby. Body: \texttt{\{code\}} or \texttt{\{lobby\_name\}}.                                           \\
        \texttt{POST}   & \texttt{/api/lobby/leave}      & Leave the current lobby.                                                                                     \\
        \texttt{GET}    & \texttt{/api/lobby/list}       & List public lobbies.                                                                                         \\
        \texttt{POST}   & \texttt{/api/lobby/broadcast}  & Broadcast a game packet to lobby peers. Body: \texttt{\{data\}}.                                             \\
        \texttt{GET}    & \texttt{/api/lobby/events}     & SSE event stream. Query: \texttt{?peer\_id=\{id\}}.                                                          \\
        \bottomrule
    \end{tabular}
\end{table}

\subsection{SSE Event Types}
\label{app:c:sse-events}

The SSE stream (\texttt{GET /api/lobby/events}) delivers the following event types:

\begin{table}[htbp]
    \centering
    \caption{SSE event types}
    \begin{tabular}{@{}lp{9cm}@{}}
        \toprule
        \textbf{Event}            & \textbf{Description}                            \\
        \midrule
        \texttt{welcome}          & Initial connection acknowledgement with peer ID \\
        \texttt{peer\_joined}     & A new peer joined the lobby                     \\
        \texttt{peer\_left}       & A peer left the lobby                           \\
        \texttt{lobby\_joined}    & Confirmation that this peer joined a lobby      \\
        \texttt{lobby\_closed}    & The lobby was closed (with reason)              \\
        \texttt{game\_packet}     & A game data packet relayed from another peer    \\
        \texttt{heartbeat}        & Keep-alive ping                                 \\
        \texttt{error}            & Error notification                              \\
        \texttt{server\_shutdown} & Server is shutting down                         \\
        \bottomrule
    \end{tabular}
\end{table}

\subsection{WebSocket Endpoints}
\label{app:c:websocket-endpoints}

These provide WebRTC signaling (for native desktop clients) and legacy WebSocket lobby handling:

\begin{table}[htbp]
    \centering
    \caption{WebSocket endpoints}
    \begin{tabular}{@{}lp{9cm}@{}}
        \toprule
        \textbf{Path}         & \textbf{Description}                                                                            \\
        \midrule
        \texttt{/ws/\{code\}} & WebRTC signaling channel. Exchanges ICE candidates, SDP offers/answers between peers in a room. \\
        \texttt{/lobby}       & WebSocket-based lobby management (native client usage; browser clients use HTTP+SSE instead).   \\
        \bottomrule
    \end{tabular}
\end{table}

\subsection{Utility Endpoints}
\label{app:c:utility-endpoints}

\begin{table}[htbp]
    \centering
    \caption{Utility endpoints}
    \begin{tabular}{@{}llp{7cm}@{}}
        \toprule
        \textbf{Method} & \textbf{Path}             & \textbf{Description}                                 \\
        \midrule
        \texttt{GET}    & \texttt{/health}          & Health check. Returns room, lobby, and peer counts.  \\
        \texttt{GET}    & \texttt{/rooms}           & List all active rooms.                               \\
        \texttt{GET}    & \texttt{/lobbies}         & List all lobbies with peer details.                  \\
        \texttt{GET}    & \texttt{/api/server/info} & Server information (version, uptime, configuration). \\
        \texttt{GET}    & \texttt{/}                & Welcome page / root info.                            \\
        \bottomrule
    \end{tabular}
\end{table}

\subsection{Error Codes}
\label{app:c:error-codes}

All error responses follow the format \texttt{\{"success": false, "error": "<code>", "message": "..."\}}.

\begin{table}[htbp]
    \centering
    \caption{Signaling server error codes}
    \begin{tabular}{@{}lp{9cm}@{}}
        \toprule
        \textbf{Code}               & \textbf{Meaning}                   \\
        \midrule
        \texttt{LOBBY\_NOT\_FOUND}  & Specified lobby does not exist     \\
        \texttt{LOBBY\_CLOSED}      & Lobby has been closed              \\
        \texttt{LOBBY\_FULL}        & Lobby has reached player limit     \\
        \texttt{ALREADY\_IN\_LOBBY} & Peer is already in a lobby         \\
        \texttt{NOT\_IN\_LOBBY}     & Peer is not in any lobby           \\
        \texttt{ROOM\_NOT\_FOUND}   & Specified room does not exist      \\
        \texttt{PEER\_NOT\_FOUND}   & Specified peer ID not found        \\
        \texttt{PEER\_ID\_IN\_USE}  & Requested peer ID is already taken \\
        \texttt{UNKNOWN\_COMMAND}   & Unrecognized message type          \\
        \texttt{INVALID\_JSON}      & Malformed JSON in request body     \\
        \bottomrule
    \end{tabular}
\end{table}

%-----------------------------------------------------------------------
\section{Server Data Models}
\label{app:c:server-models}

The signaling server uses two primary data models:

\paragraph{Peer:} Represents a connected player. Fields: \texttt{peer\_id} (int), \texttt{ws} (WebSocket connection, optional), \texttt{sse\_queue} (asyncio Queue for SSE events, optional), \texttt{player\_data} (dict), \texttt{lobby\_code} (str, optional). A peer uses either WebSocket \emph{or} SSE, not both---this dual-transport design enables both native and browser clients.

\paragraph{Lobby:} Represents a game lobby. Fields: \texttt{code} (str), \texttt{name} (str), \texttt{host\_id} (int), \texttt{public} (bool), \texttt{player\_limit} (int), \texttt{open} (bool), \texttt{created\_at} (ISO timestamp), \texttt{peers} (dict of Peer).

%-----------------------------------------------------------------------
\section{GDScript Game Classes}
\label{app:c:gdscript-classes}

This section summarizes the key GDScript classes. Full source code is available in the repository (see Appendix~\ref{app:b:repository}).

\subsection{CardManager (card\_manager.gd)}
\label{app:c:card-manager-class}

Base class managing all card logic for single-player mode.

\paragraph{Key Properties:}
\begin{itemize}
    \item \texttt{num\_cards: int} --- Number of cards in the game
    \item \texttt{cards\_array: Array} --- All card node instances
    \item \texttt{sorted\_all: Array} --- Sorted reference array (target ordering)
    \item \texttt{slots: Array} --- Buffer slot instances
    \item \texttt{card\_container: HBoxContainer} --- Shared card container (scene tree node)
\end{itemize}

\paragraph{Key Methods:}
\begin{itemize}
    \item \texttt{check\_sorting\_order() -> bool} --- Validates ascending order of cards in container
    \item \texttt{check\_player\_buffer\_contiguity() -> bool} --- Validates buffer cards form a contiguous sorted subsequence
    \item \texttt{generate\_completed\_card\_array(values, prefix) -> Array} --- Creates card instances from value array
    \item \texttt{create\_buffer\_slots() -> Array} --- Instantiates player buffer slots
    \item \texttt{animate\_card\_swap(card\_a, card\_b)} --- Animated card position swap with tween
\end{itemize}

\subsection{MultiplayerCardManager (multiplayer\_card\_manager.gd)}
\label{app:c:multiplayer-card-manager-class}

Extends \texttt{CardManager} for multiplayer mode.

\paragraph{Inner Class --- CardState:}
\begin{itemize}
    \item Serializable state object with \texttt{to\_dict()} and \texttt{from\_dict()} methods
    \item Fields: \texttt{value}, \texttt{index}, \texttt{original\_index}, \texttt{in\_container}, \texttt{in\_buffer}, \texttt{buffer\_owner}
\end{itemize}

\paragraph{Key Methods (GDSync-exposed):}
\begin{itemize}
    \item \texttt{sync\_complete\_game\_state(state\_data)} --- Full state reconciliation
    \item \texttt{sync\_card\_moved(value, from\_idx, to\_idx, client\_id)} --- Card repositioning
    \item \texttt{sync\_card\_entered\_buffer(value, player\_id)} --- Card enters public buffer
    \item \texttt{sync\_card\_left\_buffer(value, player\_id, new\_idx)} --- Card leaves buffer
    \item \texttt{barrier\_thread\_reached(player\_id)} --- Player reached barrier
    \item \texttt{barrier\_activate()} --- All threads at barrier; activate
    \item \texttt{barrier\_release()} --- Release barrier; resume normal play
\end{itemize}

\subsection{BarrierManager (barrier\_manager.gd)}
\label{app:c:barrier-manager-class}

Implements barrier synchronization state machine.

\paragraph{States:} \texttt{RUNNING} $\to$ \texttt{WAITING\_AT\_BARRIER} $\to$ \texttt{BARRIER\_ACTIVE} $\to$ \texttt{RUNNING}.

\paragraph{Key Methods:}
\begin{itemize}
    \item \texttt{enter\_waiting\_state(initiator\_id, all\_player\_ids)} --- Transition to waiting state
    \item \texttt{activate\_barrier()} --- Transition to active state (all threads arrived)
    \item \texttt{release\_barrier()} --- Reset and return to running state
    \item \texttt{determine\_main\_thread(first\_id, all\_ids) -> int} --- Select coordinator (first-to-reach or round-robin)
\end{itemize}

\paragraph{Signals:}
\begin{itemize}
    \item \texttt{barrier\_state\_changed(new\_state)} --- State transition notification
    \item \texttt{thread\_reached\_barrier(player\_id)} --- Player arrival notification
    \item \texttt{main\_thread\_assigned(player\_id)} --- Main thread selected
\end{itemize}

\subsection{ConnectionManager (connection\_manager.gd)}
\label{app:c:connection-manager-class}

Autoload that wraps GDSync signals into game-level events.

\paragraph{Key Methods:}
\begin{itemize}
    \item \texttt{am\_i\_host() -> bool} --- Whether local player is lobby host
    \item \texttt{get\_my\_client\_id() -> int} --- Local player's GDSync client ID
    \item \texttt{get\_player\_list() -> MultiplayerTypes.PlayersMap} --- Current player map
\end{itemize}

\paragraph{Key Signals:}
\begin{itemize}
    \item \texttt{player\_joined\_lobby(player\_id)} --- Player entered lobby
    \item \texttt{player\_list\_updated(players)} --- Player list changed
    \item \texttt{lobby\_joined(lobby\_name)} --- Local player joined a lobby
\end{itemize}

%-----------------------------------------------------------------------
\section{GDSync Framework Usage}
\label{app:c:gdsync-usage}

This project uses GDSync's explicit function exposure model. GDSync does \emph{not} use annotations or decorators for RPC. Instead, functions are registered at runtime:

\begin{lstlisting}[caption={GDSync function exposure (actual API)}]
# Expose a function for remote calls
GDSync.expose_func(self.my_function)

# Call function on all peers
GDSync.call_func(self.my_function, [arg1, arg2])

# Call function on a specific peer
GDSync.call_func_on(target_peer_id, self.my_function, [arg1, arg2])

# Expose a variable for synchronization
GDSync.expose_var(self, "variable_name")
\end{lstlisting}

\paragraph{Key GDSync Signals:}
\begin{itemize}
    \item \texttt{GDSync.connected} --- Successfully connected to server
    \item \texttt{GDSync.connection\_failed} --- Connection attempt failed
    \item \texttt{GDSync.lobby\_created} --- Lobby created (host)
    \item \texttt{GDSync.lobby\_joined} --- Joined an existing lobby
    \item \texttt{GDSync.client\_joined} --- A new client joined the lobby
    \item \texttt{GDSync.client\_left} --- A client left the lobby
    \item \texttt{GDSync.lobbies\_received} --- Lobby list received
\end{itemize}

\noindent For protected mode details and the web-export compatibility layer, see Chapter~\ref{ch:problems}, Section~\ref{subsec:protected-mode-issue} and Chapter~\ref{ch:implementation}, Section~\ref{sec:web-patch-implementation}.

%-----------------------------------------------------------------------
